% ============================================================
% LaTeX 环境测试文档
% 编译：XeLaTeX（推荐） / PdfLaTeX / LuaLaTeX
% 引用：Biber（推荐） / BibTeX
% 如果这份文档能编译通过，说明你的 LaTeX 环境配置成功
% ============================================================

\documentclass[a4paper,12pt]{ctexart}  % ctexart 原生支持中文，免去手动配字体

% ==== 数学 ====
\usepackage{amsmath, amssymb}

% ==== 图形与颜色 ====
\usepackage{graphicx}
\usepackage{xcolor}

% ==== 超链接 ====
\usepackage[hidelinks]{hyperref}

% ==== 代码块 ====
\usepackage{listings}
\lstset{
    basicstyle = \ttfamily\small,
    keywordstyle = \bfseries\color{blue},
    commentstyle = \color{gray},
    numbers = left,
    numberstyle = \tiny,
    frame = single,
    breaklines = true,
    showstringspaces = false
}

% ==== 参考文献（BibLaTeX + Biber 推荐）====
\usepackage[backend=biber, style=gb7714-2015]{biblatex}  % GB/T 7714 国标格式
\addbibresource{ref.bib}  % 如果还没有 ref.bib，编译时会跳过

% ============================================================

\title{\textbf{LaTeX 环境测试文档}}
\author{name}
\date{20xx年xx月xx日}

\begin{document}

\maketitle           % 标题
\tableofcontents     % 目录
\newpage

% ============================================================
\section{概述}
% ============================================================

如果你能看到这行中文，说明 \textbf{XeLaTeX}（或 LuaLaTeX）编译器和 \texttt{ctex} 宏包已经正常工作。

这份文档同时演示了以下功能：
\begin{itemize}
    \item 中英文混排、字体样式（\textbf{粗体}、\textit{斜体}、\texttt{等宽}）
    \item 数学公式（行内与行间）
    \item 图片插入
    \item 表格
    \item 代码块
    \item 交叉引用
    \item 参考文献引用
    \item 超链接
\end{itemize}

% ============================================================
\section{数学公式}
\label{sec:math}
% ============================================================

\subsection{行内与行间公式}

勾股定理：$a^2 + b^2 = c^2$，其中 $a$ 和 $b$ 是直角边，$c$ 是斜边。

欧拉公式（行为公式，无编号）：
\[
    e^{i\pi} + 1 = 0
\]

欧拉公式（带编号，方便引用）：
\begin{equation}
    e^{i\theta} = \cos\theta + i\sin\theta
    \label{eq:euler}
\end{equation}

如公式~\ref{eq:euler} 所示，欧拉恒等式是数学中最优美的等式之一。

\subsection{矩阵}

\begin{equation}
    A = \begin{bmatrix}
        a_{11} & a_{12} & \cdots & a_{1n} \\
        a_{21} & a_{22} & \cdots & a_{2n} \\
        \vdots & \vdots & \ddots & \vdots \\
        a_{m1} & a_{m2} & \cdots & a_{mn}
    \end{bmatrix}
    \label{eq:matrix}
\end{equation}

\subsection{多行公式}

\begin{align}
    \nabla \cdot \mathbf{E} &= \frac{\rho}{\varepsilon_0} \\
    \nabla \cdot \mathbf{B} &= 0 \\
    \nabla \times \mathbf{E} &= -\frac{\partial \mathbf{B}}{\partial t} \\
    \nabla \times \mathbf{B} &= \mu_0 \left( \mathbf{J} + \varepsilon_0 \frac{\partial \mathbf{E}}{\partial t} \right)
\end{align}

% ============================================================
\section{表格}
\label{sec:table}
% ============================================================

表~\ref{tab:env} 汇总了各编译引擎的适用场景。

\begin{table}[h]
    \centering
    \caption{LaTeX 编译引擎对比}
    \label{tab:env}
    \begin{tabular}{|c|c|c|c|}
        \hline
        \textbf{引擎} & \textbf{中文支持} & \textbf{字体系统} & \textbf{推荐场景} \\
        \hline
        PdfLaTeX    & 需 CJK 宏包      & Type1 / 位图   & 传统英文期刊 \\
        XeLaTeX     & 原生支持         & 系统字体        & 日常 / 中文学位论文 \\
        LuaLaTeX    & 原生支持         & 系统字体        & 复杂字体 / 多语言 \\
        \hline
    \end{tabular}
\end{table}

% ============================================================
\section{代码块}
\label{sec:code}
% ============================================================

下面是一段 C 语言的「Hello, World」示例：

\begin{lstlisting}[language=C, caption={Hello World in C}]
#include <stdio.h>

int main() {
    printf("Hello, World!\n");
    return 0;
}
\end{lstlisting}

% ============================================================
\section{参考文献}
\label{sec:ref}
% ============================================================

BibLaTeX + Biber 是现代 LaTeX 推荐的参考文献方案。

若你有 \texttt{ref.bib} 文件，可以这样引用：在文献库~\cite{example} 中……

如果还没有 \texttt{.bib} 文件也无妨——编译时会显示警告，不影响 PDF 输出。

% 打印参考文献（无 .bib 文件也不会报错，只是显示为空）
\printbibliography

% ============================================================
\section{下一步}
\label{sec:next}
% ============================================================

\begin{enumerate}
    \item 学习宏包管理：遇到排版需求时，搜索「需求名 + LaTeX 宏包」
    \item 建立自己的 \texttt{.bib} 文献库，用 Zotero / JabRef 管理
    \item 了解 Ti\textit{k}Z / PGFPlots 用于制图和数据可视化
    \item 需要模板时，去 \href{https://www.overleaf.com/latex/templates}{Overleaf Templates} 查找
\end{enumerate}

以上所有内容编译通过，你的 LaTeX 环境就完全就绪了。

\end{document}
